\documentclass[../../main.tex]{subfiles}

\begin{document}
\begin{CJK*}{UTF8}{gbsn}

\section*{Exercise 2}

Fix positive integer $K$ and
consider the $K$-nearest neighbor regression
with $\mathcal{X}$ being any set (any metric space) and $\mathcal{Y} = \mathbb{R}$
who gives a fitted model

\begin{equation*}
    \hat{f}(x; \mathcal{D})=\frac{1}{K}\sum_{i \in N(x, \mathcal{D})}y_i,
\end{equation*}
where $\mathcal{D}$ is any data set containing points from $\mathcal{X} \times \mathcal{Y}$
and $N(x, \mathcal{D})$ denotes indices of all $K$ points in $\mathcal{D}$
whose $\mathcal{X}$ component is, in some sense, closest to $x$.
We can assume that there is almost surely no ties 
by implementing any tie breaking rule.

Now let us consider the data set

\begin{equation*}
    \mathcal{D}_{n} = \curly{(X_1, Y_1), \cdots, (X_{n}, Y_{n})}.
\end{equation*}
For any $y \in \mathcal{Y}$, put

\begin{equation*}
    \mathcal{D}_{n+1}^y = \curly{(X_1, Y_1), \cdots, (X_{n}, Y_{n}), (X_{n+1}, y)}.
\end{equation*}
Prove that there exists real valued functions $a(x; \mathcal{D}_{n}, X_{n+1})$
and $b(x; \mathcal{D}_{n}, X_{n+1})$
whose values depend only on $x$, $\mathcal{D}_{n}$, and $X_{n+1}$
such that, for all $(x,y) \in \mathcal{X} \times \mathcal{Y}$:

\begin{equation*}
    \hat{f}(x; \mathcal{D}_{n+1}^y)
    = a(x; \mathcal{D}_{n}, X_{n+1}) + b(x; \mathcal{D}_{n}, X_{n+1})y.
\end{equation*}

\paragraph{Solution}
First observe the fact that the set $N(x, \mathcal{D}_{n+1}^y)$
merely reads the set of all indices $i \in \{1, \ldots, n + 1\}$
belonging to the $K$ points amongst $X_1, \ldots, X_{n+1}$
who are the closest to $x$.
Therefore, although $y$ is present in the notation,
$N(x, \mathcal{D}_{n+1}^y)$ depends only on $x$ and the points $X_1, \ldots, X_{n+1}$.
Therefore, by simply writing 

\begin{equation*}
    \hat{f}(x; \mathcal{D}_{n+1}^y)
    = \frac{1}{K}\sum_{i \in N(x, \mathcal{D}_{n+1}^y)\cap \{1, \ldots, n\}}y_i
    + \frac{1}{K}\sum_{i \in N(x, \mathcal{D}_{n+1}^y)\cap \{n+1\}}y_i
    = \frac{1}{K}\sum_{i \in N(x, \mathcal{D}_{n+1}^y)\cap \{1, \ldots, n\}}y_i
    + \frac{1}{K}\chi_{\curly{n+1 \in N(x, \mathcal{D}_{n+1}^y)}}y,
\end{equation*}
we see that it suffices to define 

\begin{equation*}
    a(x; \mathcal{D}_{n}, X_{n+1}) = \frac{1}{K}\sum_{i \in N(x, \mathcal{D}_{n+1}^y)\cap \{1, \ldots, n\}}y_i
\end{equation*}
and 
\begin{equation*}
    b(x; \mathcal{D}_{n}, X_{n+1}) = \frac{1}{K}\chi_{\curly{n+1 \in N(x, \mathcal{D}_{n+1}^y)}}.
\end{equation*}
This completes the proof.////

\end{CJK*}
\end{document}
